\chapter*{Introduction}
\addcontentsline{toc}{chapter}{Introduction}
\markboth{Introduction}{}

This report presents thirteen machine-learning experiments implemented in Python.
The algorithms are coded directly with NumPy; Pandas is used to prepare tabular
data and Matplotlib to create figures. Each experiment uses a public dataset,
reports its held-out test performance, and saves detailed output beside its
script. Each lab chapter ends with the complete Python source listing. The
scripts and prepared CSV files are in the project directory.

\section*{Experiments at a glance}
\begin{table}[H]
\centering\small
\caption{Algorithms, datasets, and prediction tasks.}
\begin{tabularx}{\linewidth}{@{}r l X l@{}}
\toprule
\textbf{Lab} & \textbf{Algorithm} & \textbf{Dataset} & \textbf{Task}\\
\midrule
1 & k-NN & Red Wine Quality (1,599; 11 features) & 3-class\\
2 & Linear and polynomial regression & Red Wine Quality (1,599) & Regression\\
3 & Gradient descent & Concrete Strength (1,030; 8 inputs) & Regression\\
4 & MLE and MAP & Diabetes Progression (442; 10 inputs) & Regression\\
5 & Multinomial Naive Bayes & SMS Spam (5,574 messages) & Text classification\\
6 & Logistic regression & Breast Cancer Diagnostic (569; 30 inputs) & Binary classification\\
7 & Perceptron & Banknote Authentication (1,372; 4 inputs) & Binary classification\\
8 & Multi-layer perceptron & Optical Digits (5,620; 64 pixels) & 10-class\\
9 & RBF network & Seeds (210; 7 inputs) & 3-class\\
10 & Support vector machine & Sonar (208; 60 inputs) & Binary classification\\
11 & Support vector regression & Auto MPG (392; 8 inputs) & Regression\\
12 & ID3 decision tree & Car Evaluation (1,728; 6 attributes) & 4-class\\
13 & Genetic algorithm & Ionosphere (351; 34 inputs) & Feature selection\\
\bottomrule
\end{tabularx}
\end{table}

For parameter selection, the scripts use cross-validation on training data or a
validation split; test data is reserved for final evaluation. Feature scaling is
fit using training data where applicable. Classification is summarized with
accuracy and class-sensitive metrics; regression uses MAE, RMSE, and $R^2$.

\section*{Running the experiments}
Create the environment and prepare the datasets from the project root:
\begin{lstlisting}[language=bash,numbers=none,frame=none]
python3 -m venv .venv
.venv/bin/python -m pip install -r requirements.txt
.venv/bin/python data/prepare_datasets.py
.venv/bin/python 01_KNN/knn.py
\end{lstlisting}
Each script writes \code{results.txt} and PDF figures into its own lab folder.
The report uses the results produced by these scripts; no machine-learning
library is used to fit the models.

\clearpage
